% Abhakseite „Das kann ich“ – Zone nur im Gesamt (\mitzone)
%% TODO Ich-kann-Sätze: Platzhalter wie in den Aufgabendateien
\begin{abhakseite}
\ifdefined\mitzone
\abhakgruppe{Kennst du schon}
\abhak{1}{Quadrieren mit dem Taschenrechner, auch Dezimalzahlen; Quadratzahlen bis zwanzig hoch zwei aus dem Kopf; Quadrat vor Strich}
\abhak{2}{Quadratwurzel mit dem Taschenrechner ziehen: erst die Summe oder Differenz unter der Wurzel ausrechnen, dann die Wurzel, dann den Näherungswert runden}
\abhak{3}{Rechtwinkliges Dreieck erkennen: Rechtwinkelmarke (Bogen mit Punkt) in der Skizze finden, Dreiecksarten}
\abhak{4}{Längeneinheiten: Ergebnis mit Einheit, cm und m vor dem Rechnen angleichen, Runden auf eine Dezimale}
\abhak{5}{Formel nach einer Größe umstellen mit Strich, hier c² = a² + b² nach a²}
\abhak{6}{Koordinaten ablesen und Differenzen bilden, auch mit negativen Zahlen}
\abhak{7}{Höhe, Seitenhöhe, Seitenkante, Mantellinie und Radius an Pyramide und Kegel benennen, Radius aus dem Durchmesser}
\abhak{8}{Eigenschaften von gleichschenkligem Dreieck, gleichschenkligem Trapez, Parallelogramm, Drachen und Rechteck: Höhe, Symmetrieachse, Diagonalen}
\abhak{9}{Strecke halbieren, halbe Differenz zweier Seiten bilden (Überstand am Trapez)}
\abhak{10}{Quadratwurzel mit dem Taschenrechner ziehen: erst die Summe oder Differenz unter der Wurzel ausrechnen, dann die Wurzel, dann den Näherungswert runden – Fehler finden}
\abhak{11}{Quadratwurzel mit dem Taschenrechner ziehen: erst die Summe oder Differenz unter der Wurzel ausrechnen, dann die Wurzel, dann den Näherungswert runden – selbst rechnen}
\fi
\abhakgruppe{Einheit 1 · Satz und Hypotenuse}
\abhak{12}{Gilt der Satz hier?}
\abhak{13}{Hypotenuse}
\abhak{14}{Quadrate über den Seiten zeichnen und Flächen vergleichen (Kästchen zählen)}
\abhak{15}{Ergebnis mit sinnvoller Genauigkeit angeben}
\abhak{16}{Hypotenuse – Fehler finden, Begründen, Anwendung, Darstellungswechsel}
\abhakgruppe{Einheit 2 · Kathete und Umkehrung}
\abhak{17}{Welche Seite ist die längste?}
\abhak{18}{Kathete}
\abhak{19}{pythagoreische Tripel erkennen und vervielfachen}
\abhak{20}{Kathete – Fehler finden, Begründen, Anwendung, Darstellungswechsel}
\abhakgruppe{Einheit 3 · Pythagoras in Figuren und Körpern}
\abhak{21}{Ganz oder halb?}
\abhak{22}{Figuren und Körper}
\abhak{23}{Rampe, Leiter, Seil an der Wand mit Skizze}
\abhak{24}{Figuren und Körper – Fehler finden, Begründen, Anwendung, Darstellungswechsel}
\end{abhakseite}
